\section{Endpoint collars of long regular critical words}
\label{sec:v45-critical-collars}
The individual edge estimate does not control a cluster of critical
values. We use a physical neighborhood whose size is independent of
word length and allow its interior margins to decrease exponentially
with distance from the two endpoints. This is possible because endpoint
influence decays faster than the permitted loss of margin.

Use arclength at the contacts, let $c_j$ be the incidence cosine, and
let $\ell_j$ be flight $j$'s length. Put
\[
 q=47/53,\qquad d_j=\min\{j,m-j\},\qquad
 e_j=q^{d_j/4},\qquad e'_j=q^{\min\{j,m-1-j\}/4}.
\]
Fix $0<\varepsilon<1/10$. A regular normal-to-normal critical word,
starting and ending in $Y_R^*$, is called $\varepsilon$-protected if
$c_j\ge\varepsilon e_j$, its collision state at time $j$ has distance
at least $\varepsilon e_j$ from $\partial Y_R^*$ in the fixed cylinder
metric, and flight $j$ has clearance at least $\varepsilon e'_j$ from
every nonincident disk. Clearance excludes the two incident disks.
The section condition applies to states outside the section as well
as to its visits. A word with a fixed margin $\varepsilon$ everywhere
is protected, but the new class also allows very small interior
incidence and decision margins. No positive source probability for
this class is presumed.

\subsection{Decay of endpoint influence}
\begin{lemma}[A uniform endpoint neighborhood]
\label{lem:v45-endpoint-collar}
There are constants $r_*>0$ and $C<\infty$, independent of
$\varepsilon,m,R$ and of the protected word, with the following
properties. Its physical reflected path extends to every endpoint pair
$|u|,|v|<\rho_\varepsilon=r_*\varepsilon^3$. On this square it retains
its center itinerary and every section-membership decision. Its exact
displacement and return index are unchanged. Internal contacts satisfy
\begin{align}
 |\partial_u y_j|&\le C\varepsilon^{-1}q^{j-1-d_j/4},&
 |\partial_v y_j|&\le C\varepsilon^{-1}q^{m-1-j-d_j/4},
                   &&1\le j\le m-1.
 \label{eq:v45-endpoint-influence}
\end{align}
Every protected margin remains at least half its stated lower bound.
\end{lemma}
\begin{proof}
At an internally reflected contact choose tangent signs successively
so that neighboring off-diagonal Hessian entries are negative. Put
$t_j=1/\ell_j$ and $D_c=\operatorname{diag}(c_1,\ldots,c_{m-1})$.
The internal action Hessian, in arclength coordinates, is exactly
$H_{\rm int}=D_c A D_c$, where
\begin{equation}\label{eq:v45-contact-matrix}
 A_{jj}=t_{j-1}+t_j+\frac{2}{Rc_j},\qquad
 A_{j,j+1}=A_{j+1,j}=-t_j.
\end{equation}
This is the second-variation identity of
Proposition~\ref{prop:general-edge}, before eliminating the internal
contacts. Write $A=D_0(I-K)$, with $D_0$ its diagonal and $K\ge0$.
Its row sums are at most
\[
 q=\frac{2(50/3)}{2(50/3)+200/47}=47/53<1.
\]
The Neumann series consequently gives
\begin{equation}\label{eq:v45-inverse-contact}
 |(A^{-1})_{ij}|\le\frac{q^{|i-j|}}{(200/47)(1-q)},\qquad
 |(H_{\rm int}^{-1})_{ij}|\le
                       \frac{Cq^{|i-j|}}{c_ic_j}
\end{equation}
for every positive internal incidence. In the endpoint derivative,
one incidence factor cancels: the endpoint--internal entry has modulus
$c_0c_1t_0$, hence
$|\partial_u y_j|\le Cc_0t_0q^{j-1}/c_j$. The other endpoint is
identical. As long as $c_j\ge\varepsilon e_j/2$, this proves
\eqref{eq:v45-endpoint-influence} and, for endpoint differentials,
\begin{align}
 |dy_j|&\le C\varepsilon^{-1}q^{3d_j/4}(|du|+|dv|),\notag\\
 |du|+|dv|+\sum_{j=1}^{m-1}|dy_j|
                &\le C\varepsilon^{-1}(|du|+|dv|).
 \label{eq:v45-summed-influence}
\end{align}
There are at most two indices at any given distance $d_j$. Empty
internal matrices, when $m=1$, require no inversion.

Continue the local implicit solution along $(su,sv)$, $0\le s\le1$.
Flight directions, incidences and collision states are uniformly
Lipschitz functions of adjacent contact positions, including at a
limiting zero incidence: their formulas divide by flight lengths, not
by incidence cosines. All lengths are bounded below by $3/50$, and
above by the fixed finite alphabet of neighboring centers. Clearance
and distance to a section boundary are Lipschitz as well. Thus the
change of a collision margin at index $j$ is at most
$C\varepsilon^{-1}q^{3d_j/4}(|u|+|v|)$; shifting to an adjacent
index changes only a fixed factor. Divide this by its allowed lower
bound $\varepsilon q^{d_j/4}$. The quotient is at most
$C\varepsilon^{-2}q^{d_j/2}(|u|+|v|)$, hence less than $1/2$ for
$|u|,|v|<r_*\varepsilon^3$ if $r_*$ is small enough. The flight
clearance estimate is the same with $e'_j$. No protected margin can
first reach half its lower bound.

For each fixed finite word the internal Hessian remains invertible.
Compactness of its contact-coordinate boxes and the strict half-margins
continue the solution to $s=1$. Choose $r_*$ also below a fixed
contact-chart radius. Local uniqueness and the endpoint injectivity
of Lemma~\ref{lem:v44-endpoint-curvature} make the continuations agree
on overlaps and supply the entire square. Positive incidence and
clearance retain physical admissibility. No collision crosses the
section boundary, so every visit, including both endpoint visits,
retains its original index. The choice of $r_*$ will be reduced once
more for relative distortion below.
\end{proof}

\subsection{Relative source distortion}
Write $F_z(u,v)$ for the reduced action in this square about critical
state $z$, and $t_z=F_z(0,0)$. Let $w_z(u,v)\,du\,dv$ be the original
section probability in these endpoint coordinates. The subscript $z$
labels a physical critical word, not a new probability law.

\begin{lemma}[Distortion relative to the intrinsic edge coefficient]
\label{lem:v45-relative-distortion}
After a fixed reduction of $r_*$, there are
$0<\lambda_*\le\Lambda_*<\infty$, independent of
$\varepsilon,m,R,z$, such that
\begin{align}
 \lambda_* I_2\le\nabla^2 F_z\le
                         \Lambda_* I_2,\qquad
 e^{-1}\le w_z(u,v)/w_z(0,0)\le e.
 \label{eq:v45-relative-collar}
\end{align}
At the critical point the coefficient of the local density jump is
\begin{equation}\label{eq:v45-intrinsic-jump}
 J_z=\frac{2\pi w_z(0,0)}{\sqrt{\det\nabla^2F_z(0,0)}}>0.
\end{equation}
It is the coefficient for $\nu_R^*$; the coefficient for $\nu$ is
$cJ_z$, where $c=91/(10000\pi)$.
\end{lemma}
\begin{proof}
The first variation is $dF_z=-p_0\,du+p_m\,dv$.
Endpoint influence bounds the changes of both endpoint tangential
velocities. They start at zero, so the endpoint cosines can be kept
above $1/2$. The reduced Hessian is an internal Schur complement.
It dominates the endpoint curvature terms and is bounded above by
the endpoint principal block of the full contact Hessian. The latter
has bounded entries, since $1/\ell_j\le50/3$. Thus one may take,
for example, a lower bound $1/(2R_+)$ and an upper bound
$2(1/R_-+50/3)$, enlarging constants if necessary. These estimates
hold throughout the convex endpoint square, not just at its center.

Let $C_z=\partial_u\partial_vF_z$. The change of variables from
$(u,p_0)$ to $(u,v)$ and the collision normalization give
\begin{equation}\label{eq:v45-cross-density}
 w_z=\frac{|C_z|}{4\pi R c},\qquad
 |C_z|=\frac{c_0c_m\prod_{j=0}^{m-1}t_j}{\det A}.
\end{equation}
The last formula follows by successive tridiagonal elimination;
$\det A=1$ for the empty matrix. The internal factors $c_j^2$ cancel
between the product of off-diagonal entries and $\det H_{\rm int}$.
In particular the formula is nonzero and may be differentiated
logarithmically on the collar. It gives
\begin{equation}\label{eq:v45-log-distortion}
 d\log|C_z|=d\log c_0+d\log c_m+
            \sum_{j=0}^{m-1}d\log t_j-\operatorname{tr}(A^{-1}dA).
\end{equation}
The differentials of $t_j$ and $c_j$ have uniformly bounded
coefficients in adjacent contact displacements. The only nonzero entries
of $dA$ are on three diagonals. Its diagonal potential derivative is
bounded by $C|dc_j|/c_j^2$, and the relevant entries of $A^{-1}$
are uniformly bounded by \eqref{eq:v45-inverse-contact}. Consequently
\eqref{eq:v45-summed-influence} bounds the potential part of the trace by
\[
 C\varepsilon^{-3}\sum_{j=1}^{m-1}q^{d_j/4}(|du|+|dv|)
       \le C\varepsilon^{-3}(|du|+|dv|).
\]
The two endpoint cosines stay above $1/2$. The remaining terms of
\eqref{eq:v45-log-distortion} are bounded by
$C\varepsilon^{-1}(|du|+|dv|)$. Hence
$|d\log w_z|\le C\varepsilon^{-3}(|du|+|dv|)$.
Integration on the square of radius $r_*\varepsilon^3$ makes the
absolute logarithmic ratio at most one after reducing $r_*$.
An absolute derivative bound for an
$m$-fold collision map has not been used in its place.

Finally the Morse lemma at the strictly convex minimum gives
\eqref{eq:v45-intrinsic-jump}, or equivalently the edge formula in
Proposition~\ref{prop:general-edge}, divided by $c$. This coefficient
is not scaled when a neighborhood is chosen.
\end{proof}

\begin{proposition}[A physical mass collar of fixed roof width]
\label{prop:v45-mass-collar}
There are $h_*>0$, $a_*>0$ and $A_*<\infty$, independent of
$\varepsilon,m,R,z$, with the following properties. Set
$h_\varepsilon=h_*\varepsilon^6$. For $0<h<h_\varepsilon$ define
\[
 U_z(h)=\{(u,v):0\le F_z(u,v)-t_z<h\}
\]
inside the endpoint square and identify it with its original source
states. This sublevel lies strictly inside the square and all its
states satisfy
\begin{equation}\label{eq:v45-normal-strip}
 |p_0|+|p_m|\le A_*\sqrt h.
\end{equation}
Its roof density $g_{z,h}$ is supported on $[t_z,t_z+h]$ and obeys
\begin{equation}\label{eq:v45-collar-density}
 a_* J_z\le g_{z,h}(t_z+s)\le A_* J_z
      \quad(0<s<h),\qquad
 \nu_R^*(U_z(h))\ge a_* hJ_z.
\end{equation}
The inequalities hold almost everywhere. Collars belonging to
different physical words are disjoint in the original source.
\end{proposition}
\begin{proof}
Strong convexity and $\nabla F_z(0)=0$ give
$\lambda_*|(u,v)|^2/2\le F_z-t_z
 \le\Lambda_*|(u,v)|^2/2$.
Choose $h_*<\lambda_*r_*^2/8$, so that
$h_\varepsilon<\lambda_*\rho_\varepsilon^2/8$.
The sublevel is then a compactly contained strictly convex disk, up to
its open outer boundary. The upper Hessian bound and the first
variation prove \eqref{eq:v45-normal-strip}.

For $0<s<h$, orient the smooth closed convex level by its normal.
Its curvature is $k=(T^{\mathsf T}\nabla^2F_zT)/|\nabla F_z|$, and
its total turning is $2\pi$. Therefore
\[
 \frac{2\pi}{\Lambda_*}\le
   \int_{\{F_z=t_z+s\}}\frac{d\mathcal H^1}{|\nabla F_z|}
       \le\frac{2\pi}{\lambda_*}.
\]
Multiply by the relative bounds for $w_z$ and use
\eqref{eq:v45-intrinsic-jump}. Coarea gives both density inequalities,
and integration in $s$ gives the mass lower bound. Distinct center
words are disjoint regular itinerary events; a single word has at
most one critical point. Thus no source mass is counted twice.
\end{proof}
